%----------------------------------------------------------------------
% 英文摘要
%----------------------------------------------------------------------

% 把英文摘要寫在裡面
\begin{enAbstract}
The spin angular momentum (SAM) of light was established in 1936, while the orbital angular momentum (OAM) carried by structured light was later demonstrated in 1992 using Laguerre--Gaussian (LG) beams. Recent studies have shown that coupling between SAM and OAM can occur in non-paraxial (strongly focused) light. However, only within the paraxial regime can a light beam carry simultaneously well-defined SAM and OAM. Previous studies have identified SAM and OAM interaction (SOI) in the phase of the small but finite longitudinal component of paraxial beams. Here, we extend these studies by superposing two LG beams to form vector vortex beams (VVBs), and show that the corresponding SOI term is preserved in the amplitude of the longitudinal field.

To provide a practical route for observing this SOI feature in VVBs, two-dimensional transition metal dichalcogenides (TMDs) have emerged as promising platforms for studying light--matter interactions, owing to their finite direct bandgap in the visible to near--infrared spectral range and their strong excitonic properties. In particular, gray excitons (GXs), which couple to light through the out--of--plane (longitudinal) field component, are well suited for this purpose. Moreover, the transition dipole moment of GXs is approximately isotropic in momentum space, allowing the resulting light--matter interaction to directly reflect the SOI of VVBs. This study finds that both valley-mixed bright excitons (BXs) and gray excitons can be selectively excited via different mechanisms by exploiting the relative phase control of VVBs.

Furthermore, we extend previous studies on the realization of valley-mixed exciton wave packets by explicitly extracting the phase structure embedded in the internal exciton wave function, such that common factors cannot be trivially removed as a global phase. While the phase term can be approximately separated for valley exciton wave packets, the situation is more involved for valley-mixed excitons. In this case, not only phase extraction is required, but also a more careful construction of the exciton wave function is necessary to avoid arbitrary phase factors that can significantly alter the resulting wave-packet profile.
    \enAbsKeywords
\end{enAbstract}
